%!TEX root = iclr2026_conference.tex
\begin{abstract}
\begin{revblock}
World models usually predict at a fixed temporal stride and state description. In action-sparse physics, one launch can trigger flight, collision and settling without further action. The lowest-error horizon need not give the best action ranking. We introduce a prediction request $r=(\Delta,\alpha)$ that advances a shared latent state by $\Delta$ frames at description level $\alpha$. One predictor serves all nine requests, and a dynamic-programming-distilled policy selects one per query. We evaluate 400 held-out, within-distribution NovPhy levels, with grid AUC on 311 mixed-outcome levels, and rollout error also on CLEVRER and Physion-Dominoes. At matched training updates per request, sharing improves nine-request mean ranking AUC over single-request specialists by $+0.0848$ [$0.0742$, $0.0952$], although specialists given more updates rank most coarse requests better. Rollout-error ordering depends on the training recipe and benchmark. On the NovPhy grid, $\Delta=1$ gives the lowest rollout position error, whereas the highest-AUC fixed request on the same predictor uses $\Delta=15$. In the pixel studies, $\Delta=15$ instead has the lower error. On one checkpoint, joint selection exceeds continuous-only selection by $+0.0201$ [$0.0154$, $0.0248$] AUC at higher compute, but does not outperform fixed ($\Delta=15$, continuous) requests on the primary continuous predictor. Direct success prediction ranks far better at the tested budgets, and adding the rollout endpoint slightly lowers AUC. Request conditioning enables controlled granularity studies. Evaluating adaptive selection requires fixed-request, no-model and direct-prediction baselines with compute reporting.
\end{revblock}
\end{abstract}
